\chapter{隔空对话}
% 完整版：chapters/05_serocomputer.tex

\pencilfig{5}{\pencilart{5}}{左边是电话：一根导线，一对一。右边是电台：一条消息同时发给周围所有人。}

到目前为止，神经元彼此之间都是打电话：导线对导线，点对点。现在打开第一座电台——释放血清素的\nt{SERO}神经元。整个大脑里它们只有几十万个，但它们的导线铺满了大脑的大片区域。

\section*{四点不同}

\nt{SERO}神经元和\nt{GLUT}不一样。第一，它的信号符号由接收者决定：血清素的“锁”两种符号都有。第二，它自己会工作：在清醒时，只要有持续的背景供给，它就以每秒几次的节奏均匀地咔嗒——这是一个节拍器。第三，它很慢：它的信号不是在几毫秒内起作用，而是在几百毫秒内。第四，它常常不把物质释放到窄缝里，而是释放到周围空间，物质像一滴墨水落进水里那样弥散开。接收者感受到的是浓度，并不知道分子是谁发来的。

\section*{并不存在的逻辑}

带抑制性“锁”的节拍器就是一个现成的“非”：没有血清素时，神经元工作；血清素一来，它就沉默。一个神经元就顶一整套电路。看起来，血清素网络的逻辑能力比\nt{GLUT}网络更强。

可是有个问题。释放到空间里的物质会混在一起。两个信号要保持不同，释放它们的位置之间的距离必须大于物质来得及弥散的范围——也就是几十微米。而每个血清素神经元都有大量释放点，散布在大片区域里，不同神经元的“导线”彼此重叠。这样的网络传不了单个的比特。它只能传少数几个全局量——就像所有人都在收听同一个台的电台。

\pencilfig{123}{%
% two release sites: far apart the clouds stay separate (two signals), close together they merge (one level)
\foreach \x/\y in {0.9/1.55, 4.1/1.55, 1.9/0, 2.9/0}{
  \foreach \r/\o in {0.95/0.07, 0.7/0.09, 0.45/0.12}{\fill[pcViolet, opacity=\o] (\x,\y) circle (\r);}
  \draw[dotpen=pcViolet, fill=pcViolet] (\x,\y) circle (0.07);}
\draw[arr] (1.95,1.55) -- (3.05,1.55); \draw[arr] (3.05,1.55) -- (1.95,1.55);
\node[lbl, anchor=south, align=center] at (2.5,1.6) {距离大于\\弥散范围};
\node[lbl, anchor=west] at (5.25,1.55) {两个信号};
\node[lbl, anchor=west] at (5.25,0) {一团共同的云};
}{血清素释放到空间里，弥散成一团云。两个来源只有相距比物质弥散范围更远时，才给出两个不同的信号。如果离得更近，两团云就融合成同一个水平。}

\section*{它实际上在做什么}

对于一个全局量，这套系统什么都不缺。血清素的浓度把一次次咔嗒抹平成一个平滑的水平——最简单的转换器就是这样把脉冲流变成电压的。这个水平能保持几秒，然后自行消退：它不是一个需要擦除的比特，而是一道慢慢消失的痕迹。

血清素神经元身上还有针对自己的“锁”：释放出的血清素会抑制它自己的来源。工程师会从中认出一个带负反馈的放大器——一个维持水平的恒温器。不过实验表明，这种自我抑制是定点作用的，只引起短暂的停顿，所以恒温器这个角色目前与其说是事实，不如说是假说。

\section*{等待的代价}

什么样的全局量值得保持缓慢、并且对整个网络都一样？一个好的候选是\emph{耐心}。设想一个选择：现在拿一份小奖励，还是几秒后拿一份大奖励。你愿意等多久，取决于未来在你眼中“贬值”得有多快。按照一种假说，转动这个旋钮的正是血清素。有实验支持它：人为激活血清素神经元，小鼠会更久地等待延迟的奖励，然后才放弃。

本章引入了三样装置，大脑的所有电台都会用到它们：节拍器神经元，其实是一整片空间的“导线”，以及用缓慢的水平代替单个的咔嗒。

\fullversion{5}
